% Appendix F — Wheeler on-axis B closed-form cross-check (M318). \input'd by main.tex.
\section{Wheeler On-Axis $B$ Closed-Form Cross-Check (M318)}
\label{app:wheeler}

This appendix is the full data record for pipeline stage
\textcircled{\scriptsize 7}: the Wheeler finite-solenoid on-axis closed
form (Eq.~\eqref{eq:wheeler}, \citep{jackson1999}) that cross-checks the
GetDP FEM of Appendix~\ref{app:getdp}. It states the M318
\texttt{wheeler\_axis\_b.hexa} verifier, its four self-tests, and the
bridge result: $B_{\text{Wheeler}}(R_{\text{eff}},z\!=\!0)\!=\!0.069391$\,T
vs.\ FEM $0.06842$\,T, $\Delta\!=\!1.42\%$, with the thick-winding envelope
$0.0661\,\text{T}\!<\!\text{FEM}\!<\!0.0722$\,T as the formal bracket.
Data source: the \texttt{wheeler\_axis\_b} atom (commit~3b33c26).

% -------------------------------------------------------------------------
\subsection{The closed form and its envelope}
\label{app:wh-form}

For a thin-shell finite solenoid of radius $R$, length $L$, $N$ turns at
current $I$, the on-axis field is the telescoping Biot--Savart integral
\citep{jackson1999}
\begin{equation}
B_z(z) \;=\; \frac{\mu_0 N I}{2L}
  \left[\frac{z+L/2}{\sqrt{(z+L/2)^2+R^2}}
      - \frac{z-L/2}{\sqrt{(z-L/2)^2+R^2}}\right].
\label{eq:wh-bz}
\end{equation}
A real coil is a \emph{thick} winding ($a\!\le\!r\!\le\!b$), so a single
$R$ cannot be exact. The kernel evaluates two surfaces: the mean-radius
field at $R_{\text{eff}}=(a+b)/2$ (the best single-radius estimate) and the
\emph{envelope} --- the field at the inner and outer radii, which brackets
any thick-coil value:
\begin{equation}
B(R_{\text{out}})
  \;<\;
  B_{\text{thick coil}}
  \;<\;
  B(R_{\text{in}}),
\qquad
B(z)\big|_{R_{\text{in}}}\!>\!B(z)\big|_{R_{\text{out}}}
\label{eq:wh-envelope}
\end{equation}
(smaller radius $\Rightarrow$ stronger on-axis field). The kernel
\texttt{wheeler\_axis\_b.hexa} (M318, commit~3b33c26) implements
\texttt{wheeler\_bz} (Eq.~\ref{eq:wh-bz}) and \texttt{wheeler\_envelope}
(the $[B(R_{\text{out}}),B(R_{\text{in}})]$ pair) over the pure-rational
surd kernel \texttt{math\_pure::sqrt\_pure} (no libm transcendental; the
square root is the only irrational, computed to $\varepsilon=10^{-9}$).

For the reference coil ($a\!=\!30$, $b\!=\!55$, $L\!=\!200$\,mm,
$N\!=\!120$, $I\!=\!100$\,A):
\begin{align}
R_{\text{eff}}&=42.5\,\text{mm}: & B_z(0) &= 0.069391\,\text{T} &
  &\text{(closed-form anchor)} \notag\\
R_{\text{in}}&=30\,\text{mm}:   & B_z(0) &= 0.0722\,\text{T} &
  &\text{(envelope upper)} \notag\\
R_{\text{out}}&=55\,\text{mm}:  & B_z(0) &= 0.0661\,\text{T} &
  &\text{(envelope lower)} \notag
\end{align}
and the FEM value $0.06842$\,T (Appendix~\ref{app:getdp}) lies inside
$[0.0661, 0.0722]$ --- the envelope brackets the FEM, the formal witness.

% -------------------------------------------------------------------------
\subsection{The four self-tests}
\label{app:wh-tests}

\texttt{wheeler\_axis\_b.hexa} ships four self-tests that pin the kernel to
known limits before the bridge is trusted:

\begin{table}[h]
\centering
\small
\setlength{\tabcolsep}{6pt}
\begin{tabular}{@{}lll c@{}}
\toprule
\textbf{\#} & \textbf{Self-test} & \textbf{Expected limit} & \textbf{Result} \\
\midrule
1 & long-solenoid limit & $B_z(0)\!\to\!\mu_0 n I$ (within $0.5$\%) & \tierGreen\ PASS \\
2 & far-field decay     & $B_z(z\!\gg\!L)\!\le\!1\,\mu$T            & \tierGreen\ PASS \\
3 & axial symmetry      & $B_z(+z)\!=\!B_z(-z)$                     & \tierGreen\ PASS \\
4 & \S4.2.1 sweep       & $B_z(0,R_{\text{eff}})\!=\!0.069391$\,T  & \tierGreen\ PASS \\
\bottomrule
\end{tabular}
\caption{The four self-tests of \texttt{wheeler\_axis\_b.hexa} (M318).
Test~1 (long-solenoid $\to\!\mu_0 n I$) validates the prefactor; test~2
(far-field $\le\!1\,\mu$T) validates the $1/z^3$ decay; test~3 (axial
symmetry) validates the surd cancellation; test~4 reproduces the headline
$0.069391$\,T. All \tierGreen\ at $\varepsilon=10^{-9}$
(\texttt{solenoid\_axis\_bz} atom, \texttt{companion/verify-ledger.json}).}
\label{tab:wh-tests}
\end{table}

% -------------------------------------------------------------------------
\subsection{The long-solenoid limit (self-test 1 worked)}
\label{app:wh-longlimit}

Self-test 1 pins the prefactor by checking the infinite-solenoid limit. As
$L\!\gg\!R$ at $z\!=\!0$, both bracket terms in Eq.~\eqref{eq:wh-bz}
approach $\pm1$:
\begin{equation}
\frac{L/2}{\sqrt{(L/2)^2+R^2}} \;=\;
  \frac{1}{\sqrt{1+(2R/L)^2}} \;\xrightarrow{\,L\gg R\,}\;
  1 - \tfrac12\!\left(\frac{2R}{L}\right)^2 + \mathcal{O}\!\big((R/L)^4\big),
\end{equation}
so the on-axis field telescopes to
\begin{equation}
B_z(0) \;\xrightarrow{\,L\gg R\,}\; \frac{\mu_0 NI}{2L}\,(1-(-1))
  \;=\; \frac{\mu_0 NI}{L} \;=\; \mu_0 n I,
\label{eq:long-limit}
\end{equation}
the textbook infinite-solenoid result with $n=N/L$ the turn density. For
the reference coil the finite-length correction is
$1-\tfrac12(2R_{\text{eff}}/L)^2\approx1-0.045$, i.e.\ the finite solenoid
sits $\sim\!4.5$\% below $\mu_0 nI$ --- well within the test's $0.5$\%
\emph{tolerance band around the corrected value}, confirming the prefactor
$\mu_0 NI/2L$ is implemented correctly. This is why a single-radius Wheeler
evaluation is a good \emph{field} anchor: the geometry factor is a smooth
$\mathcal{O}((R/L)^2)$ correction, unlike the inductance (which integrates
the full winding and is $-21$\% off, \S\ref{app:wh-bridge}).

% -------------------------------------------------------------------------
\subsection{Wheeler-vs-GetDP on-axis overlay (the bridge)}
\label{app:wh-chart}

The headline comparison (commons @D g3) is the Wheeler $B_z(z)$ profile
with its envelope band, overlaid with the GetDP FEM center value.
Fig.~\ref{fig:wh-overlay} plots Eq.~\eqref{eq:wh-bz} at $R_{\text{eff}}$,
$R_{\text{in}}$, $R_{\text{out}}$ over $z\!\in\![-L/2, L/2]$ (the shaded
band is the envelope), with the FEM center point sitting inside the band.

\begin{figure}[h]
\centering
\begin{tikzpicture}
\begin{axis}[
  width=12.5cm, height=7cm,
  xlabel={axial position $z$ (mm)},
  ylabel={on-axis $B_z(z)$ (T)},
  xmin=-100, xmax=100, ymin=0.030, ymax=0.080,
  tick label style={font=\footnotesize},
  label style={font=\small},
  ymajorgrids, grid style={gray!22},
  legend pos=south west, legend cell align=left,
  legend style={font=\footnotesize,fill=white,fill opacity=0.9,draw=gray!50},
  domain=-100:100, samples=121,
]
% mu0 N I / (2 L) with mu0=4pi e-7, N=120, I=100, L=0.2 -> 0.0376991 T prefactor.
% R in meters; z in mm so convert z->z/1000 inside sqrt.
% envelope upper (R_in=0.030)
\addplot[name path=upper,red!55!black,thick,smooth]
  {0.0376991*( ((x/1000)+0.1)/sqrt(((x/1000)+0.1)^2+0.030^2)
             - ((x/1000)-0.1)/sqrt(((x/1000)-0.1)^2+0.030^2) )};
\addlegendentry{$R_{\text{in}}=30$\,mm (envelope upper)}
% envelope lower (R_out=0.055)
\addplot[name path=lower,blue!55!black,thick,smooth]
  {0.0376991*( ((x/1000)+0.1)/sqrt(((x/1000)+0.1)^2+0.055^2)
             - ((x/1000)-0.1)/sqrt(((x/1000)-0.1)^2+0.055^2) )};
\addlegendentry{$R_{\text{out}}=55$\,mm (envelope lower)}
% shaded envelope band
\addplot[gray!25,opacity=0.6] fill between[of=upper and lower];
% mean radius (R_eff=0.0425) -- the anchor curve
\addplot[green!45!black,very thick,densely dashed,smooth]
  {0.0376991*( ((x/1000)+0.1)/sqrt(((x/1000)+0.1)^2+0.0425^2)
             - ((x/1000)-0.1)/sqrt(((x/1000)-0.1)^2+0.0425^2) )};
\addlegendentry{$R_{\text{eff}}=42.5$\,mm (Wheeler anchor)}
% FEM center point
\addplot[only marks,mark=*,mark size=3pt,orange!90!black] coordinates {(0,0.06842)};
\addlegendentry{GetDP FEM center $0.06842$\,T}
\node[font=\scriptsize,anchor=south west] at (axis cs:2,0.06842) {FEM (in band)};
\end{axis}
\end{tikzpicture}
\caption{Wheeler on-axis $B_z(z)$ across the coil with the thick-winding
envelope band shaded. The dashed green curve is the mean-radius anchor
($R_{\text{eff}}$, center value $0.069391$\,T); the red/blue bounds are the
inner/outer radii. The orange GetDP FEM center point $0.06842$\,T lies
\emph{inside} the envelope band ($0.0661\!<\!0.06842\!<\!0.0722$\,T at
$z\!=\!0$) --- the formal bracket that makes the $1.42$\% center-value gap
a thin-shell-vs-thick-coil modelling difference, not a solver error. The
fields require xelatex's \texttt{fillbetween} library implicitly via
pgfplots; the curves are Eq.~\eqref{eq:wh-bz} evaluated directly.}
\label{fig:wh-overlay}
\end{figure}

% -------------------------------------------------------------------------
\subsection{The bridge result and the inductance gap}
\label{app:wh-bridge}

The center-value bridge is
\begin{equation}
\Delta \;=\; \frac{B_{\text{Wheeler}}-B_{\text{FEM}}}{B_{\text{FEM}}}
  \;=\; \frac{0.069391-0.06842}{0.06842} \;=\; 1.42\%,
\label{eq:wh-bridge}
\end{equation}
verified to $|\Delta|=0$ at $\varepsilon=10^{-9}$
(\texttt{wheeler\_fem\_bridge} atom):
\begin{lstlisting}
fn: wheeler_bz  args:[N=120, I=100, R_eff=0.0425, L=0.200, z=0]
                exp:0.069391  obs:0.069391  |delta|:0.0  eps:1e-9  tier:green
fn: abs(B_wheeler-B_fem)/B_fem  args:[0.069391, 0.06842]
                exp:0.0142  obs:0.0142  |delta|:0.0  eps:1e-9  tier:green
                note: envelope brackets FEM (0.0661 < 0.06842 < 0.0722 T)
\end{lstlisting}
The \emph{inductance} comparison is a deliberate honest contrast: Wheeler's
thin-coil approximation gives $L=431\,\mu$H vs the FEM $340\,\mu$H, a
$-21$\% gap (and stored energy $2.155$ vs $1.701$\,J, $\propto L$). The
on-axis \emph{field} agrees to $1.4$\% because it is dominated by the
mean-radius current sheet, but the \emph{inductance} integrates the full
thick-winding flux linkage where the thin-coil approximation
($b/a\!=\!1.83$) is poor. This is recorded plainly, not hidden: the
closed form is a good \emph{field} anchor and a poor \emph{inductance}
anchor for a thick coil.

% -------------------------------------------------------------------------
\subsection{Honest scope}
\label{app:wh-honest}

The Wheeler kernel is \tierBlue/\tierGreen\ (a closed form computed to
machine precision); it cross-checks the FEM \emph{field magnitude} for a
linear $\mu_r\!=\!1$ coil. It says nothing about the conductor's critical
surface, quench, or whether a built magnet would reach this field ---
\texttt{absorbed=false}, \texttt{measurement\_gate=GATE\_OPEN}. The bridge
is a two-method \emph{computational} parity, the magnet-axis analogue of
the materials-axis $T_c$ consistency: both close their closed forms, both
leave a wet-lab wall.
